\section{Vacuum backgrounds, fixed quotients and harmonic bases}

\subsection{Background construction and descent to the quotients}

The Nil background in (BF-01) is Tanimoto\textquotesingle s homogeneous vacuum solution with $b=1$, Kasner exponents $(2/3,2/3,-1/3)$ and literature time $u^3$. Sol is the twisted Gowdy solution $P=2x,Q=0,R=t$. The biaxial Taub background uses $\ell=1,m=0$ and the conversion $\sigma=-\omega/2$. Direct substitution gives zero for the full Ricci tensor, both ADM constraints and the metric/momentum evolution residuals at the stated lapse. Time reparametrization does not change the lattice \cite{Tanimoto2004,Rendall2012,BeyerHennig2012}.

\stepcounter{equation}\begin{equation}\tag{BF-02}\label{eq:BF-02}
\begin{aligned}
\pi_{\rm Nil}^{ij}&=\mathrm{diag}(-4/(9A),-4/(9A),-u^2),\\
\pi_{\rm Sol}^{ij}&=\mathrm{diag}(-1/F,-t-1/(4t),-t-1/(4t)),\\
\pi^{ij}&=\sqrt{\det G}\,(K^{ij}-G^{ij}K),\qquad K_{ij}=\dot G_{ij}/(2N).
\end{aligned}
\end{equation}

The three Nil deck generators preserve the coframe in (QA-01). The Sol gluing preserves $(e^xdy,e^{-x}dz)$, and the $S^3$ coframe is global. Consequently the time-only $G,N$ descend to the fixed quotients. Positivity on the closed intervals follows directly from (BF-01).

\subsection{Central-nonzero Nil basis and tensor grade}

\stepcounter{equation}\begin{equation}\tag{NF-01}\label{eq:NF-01}
\begin{aligned}
\mu&=m/P_0,\quad m\in\mathbb Z\setminus\{0\},\quad 0\le r<|m|,\quad P_0=2\pi,\\
\Phi_{m,r,n}&=\sqrt{P_0}\,e^{i\mu z}\sum_{k\in\mathbb Z}
 e^{i(r+km)y}f_n\!\left(x+\frac{r+km}{\mu}\right),\\
f_n(v)&=\frac{(|\mu|/\pi)^{1/4}}{\sqrt{2^nn!}}
 H_n(\sqrt{|\mu|}v)e^{-|\mu|v^2/2}.
\end{aligned}
\end{equation}

Use normalized measure $dx\,dy\,dz/P_0^4$. The lattice shift in $x$ reindexes the sum, and the $y,z$ periods give integer phases, so (NF-01) descends to the quotient. Distinct $r,m$ are orthogonal. Unfolding the sum of $x$ intervals to the real line reduces normalization to that of $f_n$.

\stepcounter{equation}\begin{equation}\tag{NF-02}\label{eq:NF-02}
\begin{aligned}
(X-iY)\Phi_n&=\sqrt{2\mu n}\,\Phi_{n-1},&
(X+iY)\Phi_n&=-\sqrt{2\mu(n+1)}\,\Phi_{n+1},&Z\Phi_n&=i\mu\Phi_n\quad(m>0),\\
\theta^\pm&=(\sigma^1\pm i\sigma^2)/\sqrt2,&\theta^0&=\sigma^3,& L&=n+\mathrm{sign}(m)s.
\end{aligned}
\end{equation}

The symmetric tensor basis is $\theta^+\theta^+$, $(\theta^+\theta^0+\theta^0\theta^+)/\sqrt2$, $(\theta^+\theta^-+\theta^-\theta^+)/\sqrt2$, $\theta^0\theta^0$ and the two negative-spin partners, with spins $2,1,0,0,-1,-2$. For $m>0$ only components with $n=L-s\ge0$ exist. The LRS background has spin zero; ladder moves and tensor-spin moves preserve $L$. Applying the same rule to metric, momentum, lapse, shift and Cartan tensors closes a finite block. Complex conjugation gives $m<0$, with $r\mapsto-r\bmod|m|$ and $s\mapsto-s$.

The spin is a local rotation grading, not an additional global rotation KID of the quotient. Non-LRS coefficients couple $n$ to $n\pm2$ and destroy this finite-grade closure; this does not exclude other non-LRS analyses. At $m=0$ the usual base Fourier basis applies.

\subsection{Fixed Sol lattice and unfolding}

\stepcounter{equation}\begin{equation}\tag{SF-01}\label{eq:SF-01}
\begin{aligned}
\varphi&=(1+\sqrt5)/2,\quad\lambda=\varphi^2,\quad L_{\rm Sol}=\log\lambda,\\
S&=\frac1{\sqrt{1+\varphi^2}}\begin{pmatrix}1&-\varphi\\\varphi&1\end{pmatrix},\quad
D=\mathrm{diag}(\lambda^{-1},\lambda),\quad SA=DS,\\
q&=2\pi S(n_1,n_2)^T=(k_0,l_0)^T,\qquad q_j=(D^T)^jq,\\
k_0l_0&=\frac{4\pi^2}{\sqrt5}(n_1^2-n_1n_2-n_2^2)\ne0\quad(n\ne0).
\end{aligned}
\end{equation}

The fiber lattice is $S\mathbb Z^2$ with area one. The integer quadratic form in (SF-01) cannot vanish for a nonzero integer pair, so every nonzero dual orbit has both fiber components. Continuous covectors on a fiber axis are not modes of this fixed quotient.

\stepcounter{equation}\begin{equation}\tag{SF-02}\label{eq:SF-02}
\begin{aligned}
U_qf&=\sqrt{L_{\rm Sol}}\sum_{j\in\mathbb Z}f(x+jL_{\rm Sol})
 e^{i(\lambda^{-j}k_0y+\lambda^jl_0z)},\\
D_1&=\partial_\xi,\quad D_2=iK,\quad D_3=iM,
\qquad K=k_0e^{-\xi},\quad M=l_0e^\xi,\\
\mathcal D&=\{f:e^{a|\xi|}\partial_\xi^bf\in L^2(\mathbb R) \text{for all integers }a,b\ge0\}.
\end{aligned}
\end{equation}

For normalized measure $dx\,dy\,dz/L_{\rm Sol}$, $U$ is an isometry from $L^2(\mathbb R)$ onto the orbit space, by deck reindexing and fiber orthogonality. All coefficients, derivatives and formal adjoints are treated on $\mathcal D$, with zero boundary terms. Smooth compactly supported functions and Hermite functions belong to this core. This does not specify a closed operator extension or prove wellposedness.

\stepcounter{equation}\begin{equation}\tag{SF-03}\label{eq:SF-03}
[D_1,D_2]=-D_2,\quad[D_1,D_3]=D_3,\quad
e^\xi\psi_0= e^{1/4}\sum_{n=0}^{\infty}\frac{\psi_n}{2^{n/2}\sqrt{n!}}.
\end{equation}

Every coefficient in (SF-03) is nonzero, so no finite Hermite truncation is exactly closed. The infinite orbit is closed under all tensor components, constraints, gauge and evolution. Its conjugate orbit completes a real field. Fiber-zero modes instead use periodic base Fourier functions and form a distinct sector.

\subsection{Wigner blocks and reality on the three-sphere}

\stepcounter{equation}\begin{equation}\tag{TF-01}\label{eq:TF-01}
\begin{aligned}
Y^j_{mn}&=\sqrt{2j+1}\,D^j_{mn},\quad d=2j+1,\quad
\overline{Y^j_{mn}}=(-1)^{m-n}Y^j_{-m,-n},\\
D_a&=2iJ_a,\quad[D_a,D_b]=-2\epsilon_{abc}D_c,\quad
\sum_aD_a^2=-4j(j+1)I.
\end{aligned}
\end{equation}

Normalize Haar measure by the volume. Expand each of the six metric coefficients in the background coframe over all right indices. The left index supplies multiplicity; a complete real block at fixed $j$ contains $d^2$ real scalar components. Tensor coefficients obey the corresponding conjugation condition in (TF-01). We classify only $j=0,1/2,1,3/2$ explicitly and use $j=3/2$ for complete reconstruction.
